\documentclass[../principal.tex]{subfiles}
\graphicspath{{\subfix{../imagenes/}}}

\begin{document}

\ifSubfilesClassLoaded{
    \chapter{Máquinas}
}{}

\section{Osciladores}

Definiremos los osciladores en varias dimensiones, pero por sobretodo nos ocuparemos en la forma de onda que generan.

\subsection{Formas de onda}

\subsubsection{Onda sinusoidal}



\subsubsection{Onda triangular}

\subsubsection{Onda diente de sierra}

\subsubsection{Onda rectangular}

\subsubsection{Onda cuadrada}

Todas las ondas cuadradas son ondas rectangulares, pero no todas las ondas rectangulares son ondas cuadradas.

Estas ondas son utopías, ya que es imposible generar una onda cuadrada perfecta que cambia de estado de forma instantánea, pero para fines perceptuales podemos aproximarnos a esta forma de onda.

\subsection{Osciladores controlados por voltaje}

Más conocidos como VCO, por su sigla en inglés, Voltage Controlled Oscillator.

\subsection{Osciladores de frecuencia baja}

Más conocidos como LFO, por su sigla en inglés, Low Frequency Oscillator.

\newpage

\end{document}

